\chapter{Fazit und Ausblick}

<<<<<<< HEAD
Die vorliegende Arbeit zeigt, dass sich präemptive CPU-Scheduling-Verfahren in einer interaktiven Webanwendung fachlich konsistent und didaktisch wirksam aufbereiten lassen. Die zentrale Leistung liegt in der Verbindung von Simulationslogik, Visualisierung und steuerbarer Interaktion zu einem integrierten Lehrwerkzeug.

Ausgehend von den theoretischen Grundlagen wurden Anforderungen an Fachlogik, Usability und Transparenz abgeleitet und in ein umsetzbares Architekturmodell überführt. Die Implementierung realisiert diese Anforderungen in einer modularen Struktur mit klarer Trennung von Simulationskern und Darstellung. Dadurch bleiben Berechnungen nachvollziehbar, testbar und unter identischen Bedingungen reproduzierbar.

Die Anwendung erfüllt damit die Zielsetzung der Arbeit: Scheduling-Entscheidungen werden sichtbar, zeitliche Abläufe mit Kennzahlen verknüpft und Verfahren auf gemeinsamer Szenariobasis direkt vergleichbar gemacht. Der didaktische Mehrwert liegt insbesondere in der Darstellung von Entscheidungspfaden statt reiner Endergebnisse.

Die Qualitätssicherung stützt diese Einordnung durch automatisierte Tests der Kernlogik, die Prüfung zentraler Interaktionspfade sowie die Berücksichtigung grundlegender Accessibility-Anforderungen. Damit ist die Anwendung sowohl als technische Umsetzung als auch als Lerninstrument belastbar einsetzbar.

\section{Ausblick}

Für die Weiterentwicklung ergeben sich mehrere Richtungen. Erstens kann die fachliche Breite durch zusätzliche Scheduling-Verfahren und differenziertere Lastmodelle erweitert werden. Zweitens lässt sich die Analyse durch weitere Kennzahlen, Exportfunktionen und ergänzende Vergleichsmodi vertiefen. Drittens sind didaktische Erweiterungen wie geführte Lernpfade, Aufgabenserien und vorbereitete Szenariosets für den Einsatz in Lehrveranstaltungen sinnvoll.

Darüber hinaus ist eine empirische Evaluation mit Studierenden ein naheliegender nächster Schritt, um den Lerneffekt systematisch zu untersuchen. Die vorliegende Arbeit schafft dafür eine tragfähige technische und konzeptionelle Grundlage.
=======
Die vorliegende Arbeit zeigt, dass sich präemptive CPU-Scheduling-Verfahren in Form einer interaktiven Webanwendung fachlich konsistent und didaktisch nachvollziehbar vermitteln lassen. Die zentrale Leistung liegt in der Verbindung von Simulationslogik, Visualisierung und steuerbarer Interaktion zu einem zusammenhängenden Lehrwerkzeug.

Ausgehend von den theoretischen Grundlagen wurden Anforderungen an Fachlogik, Usability und Transparenz systematisch abgeleitet und in ein umsetzbares Architekturmodell überführt. Die Implementierung realisiert diese Anforderungen in einer modularen Struktur, in der Simulationskern und Darstellung klar getrennt sind. Dadurch bleiben Berechnungen nachvollziehbar, testbar und unter identischen Bedingungen reproduzierbar.

Die Anwendung erfüllt damit die Zielsetzung der Arbeit: Sie macht Scheduling-Entscheidungen sichtbar, verknüpft den zeitlichen Ablauf mit relevanten Kennzahlen und ermöglicht die direkte Gegenüberstellung mehrerer Verfahren auf gemeinsamer Szenariobasis. Der didaktische Mehrwert entsteht insbesondere dadurch, dass nicht nur Endergebnisse präsentiert werden, sondern auch die zugrunde liegenden Entscheidungspfade.

Die Qualitätssicherung stützt diese Aussage durch automatisierte Tests der Kernlogik, Prüfungen zentraler Interaktionspfade und die Berücksichtigung grundlegender Accessibility-Anforderungen. Damit ist die Anwendung sowohl als technische Umsetzung als auch als Lerninstrument belastbar einsetzbar.

\section{Ausblick}

Für eine Weiterentwicklung bieten sich mehrere Richtungen an. Erstens kann die fachliche Breite durch zusätzliche Scheduling-Verfahren und differenziertere Lastmodelle erweitert werden. Zweitens lässt sich die Analyse vertiefen, etwa durch weitere Kennzahlen, Exportfunktionen und ergänzende Vergleichsmodi. Drittens ist eine didaktische Erweiterung denkbar, beispielsweise durch geführte Lernpfade, Aufgabenserien und vorbereitete Szenariosets für den Einsatz in Lehrveranstaltungen.

Darüber hinaus wäre eine empirische Evaluation mit Studierenden ein sinnvoller nächster Schritt, um den tatsächlichen Lerneffekt systematisch zu untersuchen. Die vorliegende Arbeit schafft hierfür eine tragfähige technische und konzeptionelle Grundlage.
>>>>>>> 82e002eec67f925c0df4b4f54fa62a8404ac3abc
